\providecommand{\Res}[2]{\qty[round-mode=figures,round-precision=3]{\fpeval{#1}}{#2}}
\providecommand{\Val}[1]{\num[round-mode=figures,round-precision=4]{\fpeval{#1}}}
\providecommand{\ResSci}[2]{\qty[round-mode=figures,round-precision=3,exponent-mode=scientific]{\fpeval{#1}}{#2}}
\subsubsection*{Question a) Angle limite de résolution}
\[ \theta_{\min} = \frac{1{,}22\,\lambda}{D} = \frac{\num{1.22} \cdot \num[exponent-mode=threshold]{\vsi{lam}}}{\num[exponent-mode=threshold]{\vsi{Dp1}}} = \ResSci{1.22*\vsi{lam}/\vsi{Dp1}}{rad} \]
L'angle limite de résolution est \textbf{\ResSci{1.22*\vsi{lam}/\vsi{Dp1}}{rad}}.

\subsubsection*{Question b) Distance maximale le jour}
\begin{enumerate}[label=\arabic*)]
    \item \textbf{Variables connues et recherchées :}
    \begin{itemize}[label= -]
        \item $s = \qty{\v{ec}}{m}$, $\theta_{\min} = \ResSci{1.22*\vsi{lam}/\vsi{Dp1}}{rad}$
        \item Variable recherchée : $L$
    \end{itemize}
    \item \textbf{Équations pertinentes :} les phares sont tout juste résolus quand l'angle qu'ils sous-tendent vaut $\theta_{\min}$ : $\theta = s/L$
    \item \textbf{Isolation de la variable recherchée :}
    \[ L = \frac{s}{\theta_{\min}} = \frac{sD}{1{,}22\,\lambda} \]
    \item \textbf{Résolution numérique :}
    \[ L = \frac{\num{\v{ec}}}{\Val{1.22*\vsi{lam}/\vsi{Dp1}}} = \Res{\v{ec}*\vsi{Dp1}/(1.22*\vsi{lam})/1000}{km} \]
    \item \textbf{Réponse :} le jour, on distingue les phares jusqu'à \textbf{\Res{\v{ec}*\vsi{Dp1}/(1.22*\vsi{lam})/1000}{km}}.
\end{enumerate}

\subsubsection*{Question c) Distance maximale la nuit}
\[ L = \frac{sD}{1{,}22\,\lambda} = \frac{\num{\v{ec}} \cdot \num[exponent-mode=threshold]{\vsi{Dp2}}}{\num{1.22} \cdot \num[exponent-mode=threshold]{\vsi{lam}}} = \Res{\v{ec}*\vsi{Dp2}/(1.22*\vsi{lam})/1000}{km} \]
La nuit, on distingue les phares jusqu'à \textbf{\Res{\v{ec}*\vsi{Dp2}/(1.22*\vsi{lam})/1000}{km}} : une plus grande pupille donne une meilleure résolution.
